\section{Аналитический раздел}

\subsection{Анализ предметной области}
Предметная область проекта охватывает управление надежностью распределенных
программных систем и анализ топологии микросервисной архитектуры. Современные
системы нередко строятся на основе множества взаимодействующих компонентов, что
значительно усложняет процесс контроля их работоспособности. Система
предназначена для моделирования ИТ-инфраструктуры, что позволяет
автоматизировать поиск единых точек отказа и прогнозировать распространение
каскадных сбоев.

Основными сущностями предметной области являются компоненты и связи между ними.
Компонент может представлять собой микросервис, базу данных, брокер сообщений
или внешний программный интерфейс. Связь между компонентами показывает
зависимость одного элемента от другого и может иметь различный уровень критичности.

\subsection{Анализ существующих решений}
В сфере обеспечения отказоустойчивости распределенных систем существует
некоторое количество решений, являющихся стандартом в данной области, например
ITRS Geneos~\cite{ref:geneos} и связка Prometheus~\cite{ref:prometheus} с
Grafana~\cite{ref:grafana}.

Для изучения актуальности работы было проведено сравнение существующих решений
по следующим критериям: основное назначение, модель данных, метод выявления
проблем, анализ влияния сбоев. В результате проведенного сравнения была
составлена таблица \ref{tbl:analogs-comparison}.

\begin{center}
    \begin{tabularx}{\textwidth}{|X|X|X|}
        \caption{\label{tbl:analogs-comparison}Сравнение существующих решений} \\ \hline
        Критерий сравнения & ITRS Geneos & Prometheus + Grafana \\ \hline
        \endfirsthead

        \multicolumn{3}{l}{Продолжение таблицы \thetable{}} \\ \hline
        Критерий сравнения & ITRS Geneos & Prometheus + Grafana \\ \hline
        \endhead

        \endfoot
        \endlastfoot

        Основное назначение &
        Контроль состояния серверов и приложений в реальном времени. &
        Сбор и визуализация метрик, анализ трендов. \\ \hline

        Модель данных & Иерархическая & Временные ряды \\ \hline
        Метод выявления проблем &
        Реактивный. Сравнение текущих значений с пороговыми и правилами. &
        Реактивный. Мониторинг, отклонений, оповещения. \\ \hline

        Анализ влияния &
        Ручной / Правила. Требует сложной настройки корреляции событий. &
        Отсутствует. Показывает состояние метрик, но не моделирует зависимости
        сбоев. \\ \hline
    \end{tabularx}
\end{center}

В результате проведения анализа было выявлено, что существующие решения
позволяют проводить только реактивный анализ и не дают возможности изучать
узкие места системы на этапе проектирования. В связи с этим было решено, что
разрабатываемая система должна предоставлять возможность проактивного анализа
для выявления уязвимостей и архитектурных проблем до развертывания инфраструктуры.

\subsection{Ролевая модель}
Пользователи проектируемого приложения разделяются на несколько ролей в
зависимости от их прав доступа:
\begin{itemize}
    \item незарегистрированный пользователь --- имеет возможность только пройти
        процесс авторизации в системе или зарегистрироваться по коду приглашения;
    \item разработчик --- имеет доступ к просмотру топологии только тех систем,
        за которые отвечает его команда, и использует приложение для понимания
        последствий обновления своих сервисов;
    \item SRE-инженер --- обладает правами разработчика, а также имеет
        возможность запускать алгоритмы анализа отказоустойчивости,
        моделировать сбои и просматривать формируемые отчеты;
    \item архитектор --- обладает правами SRE-инженера, дополнительно имеет
        полные права на редактирование графа инфраструктуры, включая создание,
        изменение и удаление компонентов и связей;
    \item администратор --- управляет глобальными настройками, учетными
        записями пользователей, создает команды, генерирует приглашения и назначает роли.
\end{itemize}

На рисунке \ref{img:usecases.pdf} представлена диаграмма вариантов
использования для описанных ролей.

\jimg{usecases.pdf}{Диаграмма прецедентов}

\subsection{Формализация данных}
Для реализации заявленной функциональности по разграничению доступа и
моделированию инфраструктуры необходимо определить структуру сохраняемой
информации. Выделим сущности для управления доступом: пользователь, команда,
приглашение и система. Их описание приведено в таблице \ref{tbl:relational-entities}.

\begin{center}
    \begin{tabularx}{\textwidth}{|X|X|}
        \caption{\label{tbl:relational-entities}Сущности системы для управления
        доступом} \\ \hline
        Сущность & Информация \\ \hline
        \endfirsthead

        \multicolumn{2}{l}{{Продолжение таблицы\ \thetable{}}} \\ \hline
        Сущность & Информация \\ \hline
        \endhead

        \endfoot
        \endlastfoot

        Пользователь & Идентификатор, имя пользователя, электронная почта, хэш
        пароля, идентификатор роли, идентификатор команды \\ \hline

        Команда & Идентификатор, название команды, описание \\ \hline

        Приглашение & Идентификатор, код приглашения, электронная почта
        получателя, роль нового пользователя, срок действия, статус активности,
        идентификатор целевой команды, идентификатор активировавшего
        пользователя \\ \hline

        Система & Идентификатор, название, описание, дата создания, дата
        изменения, идентификатор команды-владельца \\ \hline
    \end{tabularx}
\end{center}

На рисунке \ref{img:er-relational.pdf} приведена диаграмма сущность-связь для
выделенных сущностей управления доступом.

Для моделирования ИТ-инфраструктуры выделены сущности, представляющие собой
элементы графа: компонент и связь. Их описание приведено в таблице
\ref{tbl:graph-entities}.

\begin{center}
    \begin{tabularx}{\textwidth}{|X|X|}
        \caption{\label{tbl:graph-entities}Сущности топологии инфраструктуры} \\ \hline
        Сущность & Информация \\ \hline
        \endfirsthead

        \multicolumn{2}{l}{{Продолжение таблицы\ \thetable{}}} \\ \hline
        Сущность & Информация \\ \hline
        \endhead

        \endfoot
        \endlastfoot

        Компонент & Идентификатор, название, тип компонента, краткое описание,
        идентификатор родительской системы \\ \hline

        Связь & Идентификатор, идентификатор исходного компонента,
        идентификатор целевого компонента, тип протокола, уровень критичности
        связи \\ \hline
    \end{tabularx}
\end{center}

На рисунке \ref{img:er-graph.pdf} представлена диаграмма сущность-связь
графовой модели данных.

\jimg{er-graph.pdf}{Диаграмма сущность-связь графовой модели данных}

\subsection{Анализ моделей баз данных}
База данных представляет собой самодокументированное собрание интегрированных
записей, структурированное согласно определенной модели. Традиционно выделяют
три этапа развития моделей данных: дореляционный, реляционный и
постреляционный~\cite{ref:db-models}.

\subsubsection{Дореляционная модель}
К дореляционным моделям относятся иерархическая и сетевая. Иерархическая модель
строится по принципу древовидной структуры, где записи связаны отношениями
предок-потомок. Реализуется связь типа один-ко-многим, при этом каждый потомок
имеет строго одного предка. Сетевая модель расширяет иерархическую, позволяя
реализовывать связи типа многие-ко-многим. В такой модели запись может иметь
несколько предков, что превращает структуру в граф. Достоинством дореляционных
моделей является высокая скорость доступа к данным, а к недостаткам относятся
значительные затраты памяти на хранение указателей и жесткость схемы.

\subsubsection{Реляционные модели}
В реляционной модели данные представляются в виде совокупности взаимосвязанных
таблиц, а манипулирование ими осуществляется с помощью языка структурированных
запросов. Основные понятия модели включают домен, атрибут, кортеж, первичный и
внешний ключи. Преимуществами реляционной модели являются высокая
согласованность данных, простота инструментария и независимость прикладного
программного обеспечения от физического способа хранения данных. Недостатками
являются снижение производительности при росте объема таблиц и жесткие
ограничения на структуру.

\subsubsection{Постреляционная модель}
Постреляционные модели призваны преодолеть ограничения реляционного подхода,
обеспечивая гибкость при работе со сложными структурами. В рамках
постреляционного подхода выделяют графовую модель, где данные представляются в
виде узлов и ребер, обладающих собственными свойствами. Такая система
оптимальна для работы с высокосвязными данными, где важна скорость обхода
сложных структур. Достоинствами постреляционных моделей являются эффективная
обработка информации со сложными связями, а основным недостатком является
сложность обеспечения строгой целостности данных по сравнению с реляционными моделями.

\subsection{Требования и ограничения к разрабатываемой базе данных}
На основе структуры выделенных данных (компонентов и связей), представляющих
собой ориентированную сеть, основным требованием к разрабатываемой базе данных
является возможность оптимизированного хранения графов и эффективного
выполнения глубокого обхода по ним. Это делает классические реляционные базы
данных неэффективными из-за медленных рекурсивных запросов.

В то же время, для обеспечения безопасности и разграничения доступа к данным
требуется спроектировать структурированное и непротиворечивое хранилище для
пользователей, команд и приглашений.

На основе поставленных требований было принято решение использовать две СУБД
для хранения данных. Для хранения топологии инфраструктуры и выполнения
алгоритмического анализа будет использоваться графовая система управления
базами данных, а для хранения учетных записей пользователей и ролей будет
примена реляционная система управления базами данных.

\subsection{Формализация задачи}
На основе проведенного анализа предметной области, существующих решений и
требований к базам данных сформулируем постановку задачи на разработку. В
рамках курсовой работы необходимо спроектировать базы данных и разработать
приложение для анализа отказоустойчивости на основе графов зависимостей.

В приложении должна быть реализована возможность регистрации и авторизации. Для
изоляции данных друг от друга в системе должна быть реализована концепция
команд. Пользователи объединяются в команды, и каждая команда имеет доступ
только к тем архитектурным схемам систем, владельцем которых она является.

Для добавления новых пользователей в команды, для администраторов должна быть
реализована возможность генерации уникальных приглашений. Использование кода
приглашения позволяет новому сотруднику зарегистрироваться и автоматически
получить необходимую роль и привязку к команде.

Основная функциональность приложения должна включать создание систем,
добавление в них компонентов и настройку связей между ними. В приложении должна
быть реализована возможность запуска анализа последствий отказа выбранного
компонента, выявления единых точек отказа и обнаружения недопустимых
циклических зависимостей для выбранной топологии системы.

\numberlesssubsection{Выводы по аналитическому разделу}
В аналитическом разделе был проведен анализ предметной области управления
надежностью распределенных систем и рассмотрены существующие решения.
Сформулированы требования к проектируемому приложению, разработана ролевая
модель пользователей и формализованы данные, подлежащие хранению. Был проведен
обзор дореляционных, реляционных и постреляционных моделей баз данных. На
основе специфики задачи было принято решение о разделении данных: применении
графовой модели для хранения архитектурных зависимостей и реляционной модели
для управления метаданными и правами доступа пользователей.
